% Libro 2 · Mundo 9 · Tablas y gráficas (app: pu-graficas)
\mundo{Tablas y gráficas}{pu-graficas}{Una gráfica es un texto hecho de formas: barras, líneas, porciones. El examen pide
leerla con precisión (¿cuánto?, ¿cuándo?, ¿qué porcentaje?), calcular medidas que resumen datos (promedio, mediana,
moda) y, sobre todo, distinguir lo que la gráfica \emph{muestra} de lo que solo parece insinuar.}{%
\item \textbf{9.1} Leer gráficas
\item \textbf{9.2} Promedio, mediana y moda
\item \textbf{9.3} Qué se puede concluir}

\leccion{pu-gr-1}{Leer gráficas}{Juego de las barras que se arrastran hasta el valor.}

\subsection{La idea: primero los ejes, después los datos}
Antes de mirar las barras, lee el \clave{título}, los \clave{ejes} y las \clave{unidades}: ¿miles o millones? ¿porcentaje
o cantidad? ¿el eje empieza en cero? Un eje que no empieza en cero exagera las diferencias: una barra que se ve el doble
de alta puede representar solo un $10\,\%$ más.

\subsubsection{Tipos de gráfica}
\begin{itemize}
  \item \clave{Barras}: comparan cantidades entre categorías.
  \item \clave{Líneas}: muestran cambios en el tiempo; la pendiente indica qué tan rápido.
  \item \clave{Circular} (torta): muestra partes de un todo; cada porción es un porcentaje y un ángulo
  ($1\,\% = 3.6^\circ$).
  \item \clave{Tablas}: dan los valores exactos; a veces conviene agregar una columna con totales.
\end{itemize}

\figura{lineas}{En una gráfica de líneas, la pendiente dice qué tan rápido cambia.}

En la gráfica de líneas, el mayor aumento está entre los meses 2 y 3 (la pendiente más empinada), aunque el valor más alto
se alcanza en el mes 4.

\begin{regla}
\begin{itemize}
  \item Lee título, ejes y unidades antes que los datos.
  \item Cambio porcentual $= \tfrac{\text{final} - \text{inicial}}{\text{inicial}} \cdot 100\,\%$.
  \item En una torta: porcentaje $\times 3.6^\circ =$ ángulo; porcentaje $\times$ total $=$ cantidad.
\end{itemize}
\end{regla}

\subsection{Ejemplos resueltos}
\begin{ejemplo}[De torta a cantidad]
En una torta, la porción de «transporte» mide $72^\circ$. Si el gasto total es \$2\,500\,000, ¿cuánto se gasta en
transporte?
\solucion $72^\circ$ es $\tfrac{72}{360} = \tfrac15 = 20\,\%$: \$500\,000.
\end{ejemplo}

\begin{ejemplo}[La mayor variación]
En la gráfica de ventas del reto, ¿entre qué meses hubo la mayor caída?
\solucion Solo hay una caída: de febrero (35) a marzo (25), 10 millones.
\end{ejemplo}

\begin{trampa}
Comparar alturas de barras sin mirar dónde empieza el eje, o confundir «aumentó en 20 millones» con «aumentó un $20\,\%$».
\end{trampa}

\begin{reto}
La gráfica muestra las ventas de una tienda (en millones de pesos). ¿En qué porcentaje aumentaron las ventas de enero a
abril?
\begin{center}
\begin{tikzpicture}
\begin{axis}[width=62mm,height=36mm,ybar,bar width=7mm,ymin=0,ymax=45,symbolic x coords={Ene,Feb,Mar,Abr},xtick=data,
  axis lines*=left,ytick={0,10,20,30,40},tick label style={font=\scriptsize,color=zink},axis line style={zink},
  nodes near coords,every node near coord/.append style={font=\scriptsize,color=zink}]
  \addplot[fill=zmoss,draw=none] coordinates {(Ene,20) (Feb,35) (Mar,25) (Abr,40)};
\end{axis}
\end{tikzpicture}
\end{center}
\opciones{$20\,\%$}{$50\,\%$}{$100\,\%$}{$200\,\%$}
\end{reto}

\begin{retoresuelto}
De 20 a 40 millones: aumentó 20 sobre 20, es decir, $100\,\%$ (las ventas se duplicaron). \textbf{Opción C.}
\end{retoresuelto}

\taller
\piso{Entrena}{leer}
\begin{ejercicios}
\item Con la gráfica de ventas: ¿cuál fue el promedio mensual? \sol{$\tfrac{120}{4} = 30$ millones.}
\item ¿Qué ángulo corresponde en una torta a una porción del $15\,\%$? \sol{$54^\circ$.}
\item En la gráfica de temperaturas, ¿en cuántos meses bajó la temperatura respecto al mes anterior? \sol{Dos (meses 5 y 6).}
\end{ejercicios}
\piso{Como en el examen}{opción múltiple}
\begin{ejercicios}
\item Con la gráfica de ventas, ¿qué fracción de las ventas de los cuatro meses corresponde a febrero?
\opciones{$\tfrac{7}{24}$}{$\tfrac{35}{100}$}{$\tfrac14$}{$\tfrac13$}
\sol{A. $\tfrac{35}{120}$.}
\end{ejercicios}

\bitacora

\leccion{pu-gr-2}{Promedio, mediana y moda}{Juego de la balanza de datos.}

\subsection{La idea: tres maneras de decir «lo típico»}
\begin{itemize}
  \item La \clave{media} (promedio) reparte todo por igual: suma $\div$ cantidad. Es sensible a los valores extremos.
  \item La \clave{mediana} es el valor del medio cuando los datos se ordenan (con una cantidad par, el promedio de los dos
  del medio). Los extremos no la mueven.
  \item La \clave{moda} es el valor que más se repite. Puede haber varias o ninguna.
\end{itemize}

\subsubsection{Un caso que lo muestra todo}
Salarios de una pequeña empresa (en millones): 1.5, 1.5, 1.5, 2, 2.5 y 12 (el del dueño). Media: $\tfrac{21}{6} = 3.5$;
mediana: $\tfrac{1.5 + 2}{2} = 1.75$; moda: 1.5. Decir «el salario promedio es 3.5 millones» es cierto pero engañoso:
cinco de seis ganan menos.

\subsubsection{Rango}
El \clave{rango} es la diferencia entre el mayor y el menor; mide qué tan dispersos están los datos.

\figura{mediana}{Seis salarios: cinco bajos y uno muy alto.}

\begin{regla}
\begin{itemize}
  \item Media $= \tfrac{\sum x}{n}$. Mediana: ordenar primero. Moda: el más frecuente.
  \item Con datos en tabla de frecuencias, la media es $\tfrac{\sum x \cdot f}{\sum f}$.
  \item Un valor extremo mueve la media mucho más que la mediana.
\end{itemize}
\end{regla}

\subsection{Ejemplos resueltos}
\begin{ejemplo}[Tabla de frecuencias]
En una encuesta sobre hermanos: 0 hermanos, 4 personas; 1, 8; 2, 6; 3, 2. Halla media, mediana y moda.
\solucion Media: $\tfrac{0 + 8 + 12 + 6}{20} = 1.3$. Mediana: con 20 datos, los lugares 10 y 11; los primeros 4 son 0 y
los siguientes 8 son 1 (lugares 5 a 12): mediana 1. Moda: 1.
\end{ejemplo}

\begin{ejemplo}[Cambiar un dato]
Si en 2, 4, 6, 8, 10 el 10 se cambia por 50, ¿qué pasa con la media y la mediana?
\solucion La media pasa de 6 a 14; la mediana sigue en 6.
\end{ejemplo}

\begin{trampa}
Calcular la mediana sin ordenar los datos.
\end{trampa}

\begin{reto}
Los datos son: 3, 7, 7, 2, 9, 4, 7, 5. ¿Cuál afirmación es correcta?
\opcionesd{media $<$ mediana $<$ moda}{moda $<$ media $<$ mediana}{mediana $<$ media $<$ moda}{las tres son iguales}
\end{reto}

\begin{retoresuelto}
Media: $\tfrac{44}{8} = 5.5$. Ordenados: 2, 3, 4, 5, 7, 7, 7, 9: mediana $\tfrac{5 + 7}{2} = 6$. Moda: 7. Entonces
$5.5 < 6 < 7$. \textbf{Opción A.}
\end{retoresuelto}

\taller
\piso{Entrena}{calcular}
\begin{ejercicios}
\item Media, mediana y moda de 5, 1, 4, 4, 6. \sol{4, 4 y 4.}
\item Mediana de 12, 3, 8, 15, 9, 10. \sol{Ordenados: 3, 8, 9, 10, 12, 15: $\tfrac{9 + 10}{2} = 9.5$.}
\item ¿Qué número hay que agregar a 4, 6, 9 para que la media sea 7? \sol{$28 - 19 = 9$.}
\end{ejercicios}
\piso{Como en el examen}{opción múltiple}
\begin{ejercicios}
\item Cinco números tienen media 10 y mediana 12. Si el mayor aumenta 5, entonces
\opcionesd{la media sube 1 y la mediana no cambia}{las dos suben 1}{la media sube 5}{la mediana sube 5}
\sol{A. La suma sube 5: la media sube $\tfrac55 = 1$; el del medio no cambia.}
\end{ejercicios}

\bitacora

\leccion{pu-gr-3}{Qué se puede concluir}{Juego de las afirmaciones que la gráfica sostiene o no.}

\subsection{La idea: lo que dicen los datos y lo que no}
Las preguntas de «se puede concluir» tienen opciones de cuatro tipos:
\begin{enumerate}
  \item \clave{Correctas}: se calculan o se leen directamente de los datos.
  \item \clave{Falsas}: contradicen los datos (un cálculo mal hecho, una lectura equivocada).
  \item \clave{No se sabe}: pueden ser ciertas, pero los datos no lo dicen (causas, opiniones, datos de otras personas).
  \item \clave{Exageradas}: generalizan («todos», «nunca») a partir de una muestra.
\end{enumerate}
Solo el tipo 1 es la respuesta. El error típico es elegir una opción del tipo 3 porque suena razonable.

\subsubsection{Correlación no es causa}
Si en una gráfica los barrios con más parques tienen menos estrés, no se puede concluir que los parques \emph{causen}
menos estrés: puede haber otra explicación (más ingresos, por ejemplo). Los datos muestran que van juntos, no por qué.

\begin{regla}
\begin{itemize}
  \item ¿Se puede calcular o leer de la gráfica? Entonces se puede concluir.
  \item ¿Habla de causas, gustos, intenciones o de otras poblaciones? Probablemente no se puede.
  \item Desconfía de «todos», «nadie», «siempre», «el mejor».
\end{itemize}
\end{regla}

\subsection{Ejemplos resueltos}
\begin{ejemplo}[Una afirmación que no se sabe]
Con los datos del reto: «A los hombres les gustan más las Matemáticas que a las mujeres». ¿Se puede concluir?
\solucion No: los datos no separan por sexo.
\end{ejemplo}

\begin{ejemplo}[Una afirmación calculable]
«Los que prefieren Sociales o Arte son la mitad de los que prefieren Matemáticas o Lenguaje.» ¿Es cierta?
\solucion Sociales o Arte: $25\,\%$; Matemáticas o Lenguaje: $55\,\%$. La mitad de 55 es 27.5, no 25: es falsa.
\end{ejemplo}

\begin{trampa}
Escoger lo que uno cree cierto en la vida real en vez de lo que los datos permiten afirmar.
\end{trampa}

\begin{reto}
Se preguntó a 200 estudiantes por su materia favorita: Matemáticas $30\,\%$, Lenguaje $25\,\%$, Ciencias $20\,\%$,
Sociales $15\,\%$ y Arte $10\,\%$. ¿Cuál afirmación se puede concluir de los datos?
\opcionesl{60 estudiantes prefieren Matemáticas.}{Matemáticas es la materia más fácil.}{Más de la mitad prefiere Ciencias o Lenguaje.}{A ningún estudiante le gusta Arte.}
\end{reto}

\begin{retoresuelto}
El $30\,\%$ de 200 es 60: se calcula directamente. B es una opinión, C es falsa ($45\,\%$ no es más de la mitad) y D es
falsa ($10\,\%$ prefiere Arte, y además «preferir» no es lo mismo que «gustar»). \textbf{Opción A.}
\end{retoresuelto}

\taller
\piso{Entrena}{concluir}
\begin{ejercicios}
\item Con los datos del reto, ¿cuántos estudiantes más prefieren Lenguaje que Sociales? \sol{$10\,\%$ de 200: 20.}
\item ¿Se puede concluir que los 200 estudiantes son de un mismo colegio? \sol{No: los datos no lo dicen.}
\end{ejercicios}
\piso{Como en el examen}{opción múltiple}
\begin{ejercicios}
\item Una gráfica muestra que en una ciudad las ventas de helado y los ahogamientos aumentan en los mismos meses. Se puede concluir que
\opcionesl{comer helado causa ahogamientos.}{las dos cosas tienden a aumentar juntas en esa ciudad.}{prohibir el helado reduciría los ahogamientos.}{en todas las ciudades pasa lo mismo.}
\sol{B. Van juntas (probablemente por el calor), pero eso no muestra causa ni se generaliza a otras ciudades.}
\end{ejercicios}

\bitacora

\evaluacion
\begin{ejercicios}
\item En una torta, una porción de $90^\circ$ representa
\opciones{$25\,\%$}{$90\,\%$}{$9\,\%$}{$30\,\%$}
\sol{A. (9.1)}
\item Una gráfica de líneas muestra que una población pasó de 800 a 1\,000 habitantes. El aumento porcentual fue
\opciones{$20\,\%$}{$25\,\%$}{$200\,\%$}{$80\,\%$}
\sol{B. $\tfrac{200}{800}$. (9.1)}
\item La mediana de 7, 2, 9, 4, 6 es
\opciones{4}{6}{5.6}{7}
\sol{B. Ordenados: 2, 4, 6, 7, 9. (9.2)}
\item Las notas de un curso son: 3.0 (5 estudiantes), 4.0 (10) y 5.0 (5). La media es
\opciones{4.0}{3.5}{4.5}{12}
\sol{A. $\tfrac{15 + 40 + 25}{20}$. (9.2)}
\item Una encuesta en un barrio muestra que el $70\,\%$ de los encuestados usa bus. Se puede concluir que
\opcionesl{el $70\,\%$ de los encuestados de ese barrio usa bus.}{el $70\,\%$ de la ciudad usa bus.}{el bus es el mejor transporte.}{nadie usa bicicleta.}
\sol{A. (9.3)}
\end{ejercicios}
